% !TEX program = xelatex
% 공통수학2 · Ⅱ. 집합과 명제 · 2. 명제 · 02 명제 사이의 관계 · 1/2차시
% 학생용 활동지 (답안 없음)
\documentclass[11pt,a4paper]{article}
\usepackage{kotex}
\usepackage{iftex}
\ifPDFTeX\else
  \setmainhangulfont{Noto Serif CJK KR}
  \setsanshangulfont{Noto Sans CJK KR}
\fi
\usepackage[margin=2.1cm]{geometry}
\usepackage{parskip}
\usepackage{enumitem}
\usepackage{array}
\usepackage{tabularx}
\usepackage{xcolor}
\usepackage{amssymb}
\usepackage{tikz}
\usepackage[most]{tcolorbox}
\usepackage{fancyhdr}
\usepackage{titlesec}

\definecolor{goldc}{HTML}{B07C22}
\definecolor{goldstrong}{HTML}{8A5F16}
\definecolor{indigoc}{HTML}{2B3B57}
\definecolor{indigosoft}{HTML}{6E82A6}
\definecolor{linec}{HTML}{D6DACB}
\definecolor{surface2}{HTML}{E4E8DC}
\definecolor{writeline}{HTML}{C7CCBA}
\definecolor{inksoft}{HTML}{52604F}

\titleformat{\section}{\normalfont\sffamily\bfseries\large\color{indigoc}}{\thesection}{0.6em}{}
\titlespacing{\section}{0pt}{14pt}{6pt}
\newtcolorbox{goalbox}{enhanced, breakable, colback=white, colframe=goldc, boxrule=0.5pt, leftrule=3pt, arc=2pt, left=10pt, right=10pt, top=8pt, bottom=8pt}
\newtcolorbox{sheetbox}{enhanced, breakable, colback=white, colframe=linec, boxrule=0.6pt, arc=3pt, left=11pt, right=11pt, top=9pt, bottom=9pt}
\newcommand{\writespace}[1]{%
  \par\nobreak\vspace{3pt}
  \foreach \n in {1,...,#1} {%
    \noindent\textcolor{writeline}{\rule{\linewidth}{0.4pt}}\par\vspace{0.62cm}%
  }
}
\newcommand{\fillblank}[1]{\makebox[#1]{\hrulefill}}
\pagestyle{fancy}\fancyhf{}\renewcommand{\headrulewidth}{0pt}
\fancyfoot[L]{\footnotesize\color{inksoft} 공통수학2 · 2026학년도 2학기 · 지평선고등학교}
\fancyfoot[R]{\footnotesize\color{inksoft} 다음 시간 --- 02 명제 사이의 관계(2)}

\begin{document}

\noindent
\begin{minipage}[b]{0.62\linewidth}
  {\sffamily\footnotesize\color{indigosoft} 공통수학2 · Ⅱ. 집합과 명제 · 2. 명제}\\[2pt]
  {\sffamily\bfseries\Large 02 명제 사이의 관계 \textperiodcentered\ 28차시}
\end{minipage}\hfill
\begin{minipage}[b]{0.36\linewidth}
  \raggedleft\small\color{inksoft}
  반\ \fillblank{1.6cm} \quad 번호\ \fillblank{1.6cm}\\[4pt]
  이름\ \fillblank{3.6cm}
\end{minipage}

\vspace{4pt}
\noindent{\color{goldc}\rule{\linewidth}{2.2pt}}
\vspace{10pt}

\begin{goalbox}
\textbf{오늘의 목표}\\
명제 $p\to q$의 참·거짓을 판별하고, 그 역과 대우를 구하여 원래 명제와 대우의 참·거짓이 항상 일치함을 설명할 수 있다.
\vspace{4pt}
\small\color{inksoft}
$\square$ 자신 있음 \quad $\square$ 조금 헷갈림 \quad $\square$ 처음 봄 \hfill (수업 전 나의 상태에 표시)
\end{goalbox}

\section{생각 열기 --- 참일까 거짓일까}

\begin{sheetbox}
다음 두 명제의 참·거짓을 예측해 보자.

\begin{enumerate}[leftmargin=1.4em, itemsep=2pt]
  \item $x\le6$이면 $x^2\le36$이다.
  \item 4의 약수는 8의 약수이다.
\end{enumerate}
\writespace{3}
\end{sheetbox}

\section{개념 정리 --- 명제의 참·거짓, 역과 대우}

\begin{sheetbox}
두 조건 $p$, $q$의 진리집합을 각각 $P$, $Q$라 할 때, 명제 $p\to q$는 $P\ \fillblank{1.6cm}\ Q$이면 참이고, 그렇지 않으면 거짓이다.

\vspace{6pt}
\textbf{반례} --- 명제 $p\to q$가 거짓임을 보이려면 $p$는 참이지만 $q$는 \fillblank{2.8cm}인 예를 하나 찾으면 된다. (예: `소수는 홀수이다'의 반례는 $x=\ \fillblank{1.2cm}$)

\vspace{8pt}
\textbf{역과 대우} --- 명제 $p\to q$에 대하여
\begin{itemize}[leftmargin=1.6em, itemsep=4pt]
  \item $q\to p$를 원래 명제의 \fillblank{2.4cm}이라 한다.
  \item $\sim q\to\ \fillblank{1.2cm}$를 원래 명제의 대우라 한다.
\end{itemize}

\vspace{6pt}
\textbf{보기} --- `$x=1$이면 $x^2=1$이다.'\\
역: \fillblank{7cm}\\
대우: \fillblank{7cm}
\end{sheetbox}

\section{연습 문제}

\begin{sheetbox}
\textbf{1.} 다음 명제의 역과 대우를 구하고, 각각의 참·거짓을 판별하시오.\\
\hspace*{1.6em}⑴ $a=b$이면 $a^2=b^2$이다.\\
\hspace*{1.6em}⑵ $x^2>4$이면 $x>2$이다.
\writespace{6}

\vspace{6pt}
\textbf{2.} 대우를 이용하여 다음 명제가 참임을 증명하시오.\\
\hspace*{1.6em}`$3a+7$이 홀수이면 $a$는 짝수이다.' (단, $a$는 자연수)
\writespace{5}
\end{sheetbox}

\section{사고의 가시화 --- Connect · Extend · Challenge}

\noindent{\footnotesize\color{inksoft} 오늘 배운 `명제의 역과 대우'를 떠올리며 아래 세 칸을 채워 보자.}\\[6pt]
\begin{tabularx}{\linewidth}{@{}XXX@{}}
\textbf{\footnotesize\color{indigoc} Connect --- 무엇과 연결되나?} &
\textbf{\footnotesize\color{indigoc} Extend --- 어디까지 확장됐나?} &
\textbf{\footnotesize\color{indigoc} Challenge --- 아직 궁금한 것은?} \\[4pt]
\end{tabularx}
\noindent
\begin{minipage}[t]{0.32\linewidth}\writespace{3}\end{minipage}\hfill
\begin{minipage}[t]{0.32\linewidth}\writespace{3}\end{minipage}\hfill
\begin{minipage}[t]{0.32\linewidth}\writespace{3}\end{minipage}

\section{마무리 성찰}

\begin{sheetbox}
\begin{minipage}[t]{0.48\linewidth}
{\footnotesize\color{inksoft} 오늘 배운 것을 한 문장으로}
\writespace{2}
\end{minipage}\hfill
\begin{minipage}[t]{0.48\linewidth}
{\footnotesize\color{inksoft} 감사 일기 / 유념 한 줄}
\writespace{2}
\end{minipage}
\end{sheetbox}

\end{document}
